\chapter{Work Plan}
\label{chapter:workplan}

This chapter presents the planned development trajectory for the remainder of
the dissertation, covering the tasks to be completed, a schedule with key
milestones, and an assessment of risks with corresponding mitigation strategies.

% ----------------------------------------------------------------------
\section{Planned Tasks}
\label{section:tasks}

The remaining work builds on the functional prototype, validated firmware,
and preliminary canine dataset described in Chapter~\ref{chapter:preliminary}.
It is structured into four phases.

\textbf{Phase~2 — Sensor refinement and structured data collection.}
The first priority is to integrate the IMU mounting calibration into the main
firmware, aligning the sensor frame with the prosthetic anatomical frame.
This is followed by structured acquisition campaigns with canine subjects
fitted with the prosthesis, conducted across multiple sessions and days to
capture variability in adaptation state and gait conditions.
Sessions will be conducted in coordination with a veterinary clinic or
rehabilitation centre.
Data quality control and manual annotation of key gait events will be
performed concurrently with acquisition.
A preliminary gait segmentation algorithm — replacing the threshold-based
heuristic currently in firmware — will be developed and validated against
the annotated sessions.

\textbf{Phase~3 — Signal processing and feature extraction.}
Raw CSV signals will be processed to extract the biomechanical features
described in Chapter~5: force-related metrics
(peak load, load consistency), temporal metrics (stride duration, step
regularity, phase timing), and signal variability metrics (RMS variability,
inter-stride coefficient of variation).
Feature extraction will be validated on the annotated dataset before
proceeding to model development.

\textbf{Phase~4 — Machine learning model development.}
The extracted feature vectors will be used to train and evaluate supervised
classification models for prosthesis adaptation state assessment.
A Random Forest classifier will be used as the primary model, with logistic
regression and support vector machines as baselines.
Model selection will be based on cross-validated performance
(accuracy, precision, recall, F1-score) and interpretability of the learned
feature importances.
Leave-one-session-out cross-validation and data augmentation will be used to
address the limited dataset size expected at this stage.

\textbf{Phase~5 — Validation and dissertation writing.}
The complete system will be validated end-to-end, from data acquisition to
adaptation state classification.
Longitudinal consistency of the predictions across multiple sessions will be
assessed to confirm that the model captures genuine adaptation trends.
The dissertation document will be completed concurrently, integrating all
experimental results and analysis.

% ----------------------------------------------------------------------
\section{Milestones}
\label{section:milestones}

Table~\ref{tab:milestones} summarises the four project milestones, their
target dates, and the associated deliverables.

\begin{table}[h]
  \centering
  \caption{Project milestones and deliverables.}
  \label{tab:milestones}
  \begin{tabular}{cp{3.8cm}p{2.8cm}p{5.0cm}}
    \toprule
    \textbf{ID} & \textbf{Milestone} & \textbf{Target date} & \textbf{Deliverable} \\
    \midrule
    M1 & Prototype + PIC2 report
       & 06/2026 \textit{(completed)}
       & Validated hardware, firmware, preliminary dataset, PIC2 report \\
    M2 & Structured dataset complete
       & 09/2026
       & Annotated multi-session canine locomotion dataset \\
    M3 & ML model validated
       & 11/2026
       & Trained classifier, performance evaluation report \\
    M4 & Dissertation submitted
       & 05/2027
       & Final dissertation document \\
    \bottomrule
  \end{tabular}
\end{table}

% ----------------------------------------------------------------------
\section{Schedule}
\label{section:schedule}

Figure~\ref{fig:gantt} presents the project schedule as a Gantt chart,
covering the full dissertation period from January~2026 to February~2027.
Completed tasks are shown in dark grey; planned tasks are shown in blue.
Milestone markers are shown in orange.
The dashed red line marks the current date (June~2026).

\begin{figure}[p]
  \centering
  \begin{ganttchart}[
      hgrid,
      vgrid={*{1}{dotted}},
      x unit        = 0.62cm,
      y unit title  = 0.60cm,
      y unit chart  = 0.65cm,
      title height  = 1,
      title label font      = \scriptsize\bfseries,
      bar label font        = \small,
      group label font      = \small\bfseries,
      milestone label font  = \small\itshape,
      bar/.style            = {fill=blue!45, rounded corners=2pt},
      bar incomplete/.style = {fill=gray!55, rounded corners=2pt},
      group/.style          = {fill=gray!25, draw=black},
      milestone/.style      = {fill=orange!80, draw=orange!60!black,
                               shape=diamond, inner sep=2pt},
      link/.style           = {->, thick, gray!70},
      today        = 6,
      today rule/.style     = {red, thick, dashed},
      today label  = {Now},
      progress label text   = {},
    ]{1}{17}

    %% Title row: year labels
    \gantttitle{2026}{12}
    \gantttitle{2027}{5} \\

    %% Month labels (numbers)
    \gantttitle{1}{1}\gantttitle{2}{1}\gantttitle{3}{1}%
    \gantttitle{4}{1}\gantttitle{5}{1}\gantttitle{6}{1}%
    \gantttitle{7}{1}\gantttitle{8}{1}\gantttitle{9}{1}%
    \gantttitle{10}{1}\gantttitle{11}{1}\gantttitle{12}{1}\\
    \gantttitle{1}{1}\gantttitle{2}{1}\gantttitle{3}{1}%
    \gantttitle{4}{1}\gantttitle{5}{1}\\

    %% ── Phase 1 (complete) ──────────────────────────────────────────
    \ganttgroup{Phase 1 — Prototype \& PIC2}{1}{6} \\
    \ganttbar[progress=100, bar/.style={fill=gray!55, rounded corners=2pt}]
      {Hardware assembly}{1}{3} \\
    \ganttbar[progress=100, bar/.style={fill=gray!55, rounded corners=2pt}]
      {Firmware development}{2}{5} \\
    \ganttbar[progress=100, bar/.style={fill=gray!55, rounded corners=2pt}]
      {Sensor validation + calibration}{4}{5} \\
    \ganttbar[progress=100, bar/.style={fill=gray!55, rounded corners=2pt}]
      {Preliminary data acquisition}{5}{6} \\
    \ganttbar[progress=95, bar/.style={fill=gray!45, rounded corners=2pt}]
      {PIC2 report}{4}{6} \\
    \ganttmilestone{M1 — Prototype + PIC2 complete}{6} \\[grid]

    %% ── Phase 2 ─────────────────────────────────────────────────────
    \ganttgroup{Phase 2 — Data Collection}{6}{9} \\
    \ganttbar{IMU calibration integration}{6}{7} \\
    \ganttbar{Hosp. Vet. Lisboa + session planning}{6}{7} \\
    \ganttbar{Structured data collection sessions}{7}{9} \\
    \ganttbar{Data annotation and quality control}{7}{9} \\
    \ganttbar{Gait segmentation algorithm}{8}{9} \\
    \ganttmilestone{M2 — Dataset complete}{9} \\[grid]

    %% ── Phase 3 ─────────────────────────────────────────────────────
    \ganttgroup{Phase 3 — Signal Processing \& ML}{9}{11} \\
    \ganttbar{Feature extraction pipeline}{9}{10} \\
    \ganttbar{ML model training and evaluation}{10}{11} \\
    \ganttmilestone{M3 — Model validated}{11} \\[grid]

    %% ── Phase 4 ─────────────────────────────────────────────────────
    \ganttgroup{Phase 4 — Validation \& Writing}{11}{17} \\
    \ganttbar{End-to-end system validation}{11}{13} \\
    \ganttbar{Dissertation writing}{10}{17} \\
    \ganttmilestone{M4 — Dissertation submitted}{17} \\

    %% ── Dependencies ────────────────────────────────────────────────
    \ganttlink{elem5}{elem7}
    \ganttlink{elem10}{elem13}
    \ganttlink{elem13}{elem16}

  \end{ganttchart}
  \caption{Project schedule (01/2026 -- 05/2027).
           The dashed red line marks the current date (06/2026).
           Dark grey bars represent completed tasks; blue bars represent
           planned work. Orange diamonds mark milestones.}
  \label{fig:gantt}
\end{figure}

% ----------------------------------------------------------------------
\section{Risks and Contingency Plan}
\label{section:risks}

Table~\ref{tab:risks} identifies the main risks to the project schedule and
quality, together with their assessed likelihood and impact, and the
corresponding mitigation strategies.

\begin{table}[ht]
  \centering
  \caption{Risk assessment and mitigation strategies.}
  \label{tab:risks}
  \begin{tabular}{p{3.6cm} c c p{6.0cm}}
    \toprule
    \textbf{Risk} & \textbf{Likelihood} & \textbf{Impact} & \textbf{Mitigation} \\
    \midrule
    Difficulty obtaining canine subjects for structured sessions
    & Low & High
    & An ongoing collaboration with the Hospital Veterinário de Lisboa
      (São Bento) has been established, providing access to real canine
      patients under clinical follow-up. Sessions are coordinated with the
      hospital's clinical team, significantly reducing the risk of subject
      unavailability. Additional sessions with the current subject serve as
      a fallback if scheduling constraints arise. \\
    \addlinespace
    Insufficient labelled data for ML training
    & Medium & High
    & Apply leave-one-session-out cross-validation; explore data
      augmentation (time warping, noise injection); consider
      transfer learning from related datasets as a prior. \\
    \addlinespace
    Sensor displacement between sessions
    & Medium & Medium
    & Design a mechanical fixture for repeatable sensor positioning;
      apply session-level normalisation to reduce sensitivity to
      absolute orientation offset. \\
    \addlinespace
    IMU axis misalignment affecting feature quality
    & Low & Medium
    & IMU mounting calibration matrix has been derived;
      integration into firmware is planned as the first task
      of Phase~2, before data collection begins. \\
    \addlinespace
    Battery autonomy insufficient for field sessions
    & Low & Low
    & Wi-Fi-off mode is already implemented, reducing current draw
      to $\approx$45\,mA and extending autonomy beyond 12\,h;
      higher-capacity cells can replace the current 500\,mAh pack
      if needed. \\
    \bottomrule
  \end{tabular}
\end{table}
